\bookchapter{Bab 22: Bagaimana AI bekerja di ponsel}

\chapterintro{1. Kamera AI: “Ruang Gelap Pasca-Produksi” Anda Pindah ke Ponsel}

\booksection{Kamera AI: “Ruang Gelap Pasca-Produksi” Anda Pindah ke Ponsel}

2. Asisten Suara: Bagaimana “Sekretaris Suara” Anda Memahami Anda?

\booksection{Asisten Suara: Bagaimana “Sekretaris Suara” Anda Memahami Anda?}

\begin{enumerate}
  \item \textbf{: Mengubah gelombang suara menjadi teks.}Memahami (pemahaman bahasa alami)
  \item \textbf{: Memahami maksud, mengekstrak entitas kunci.}Menjalankan dan menjawab
  \item \textbf{: Memanggil App atau layanan yang sesuai, lalu "berbicara" kepada Anda melalui sintesis suara.}3. AI Sisi Perangkat vs. AI Cloud
\end{enumerate}

\booksection{AI Sisi Perangkat vs. AI Cloud}

4. Model Besar AI dan Ponsel: Asisten Cerdas Pribadi Anda

\booksection{Model Besar AI dan Ponsel: Asisten Cerdas Pribadi Anda}

2025-2026, ledakan large model AI (seperti DeepSeek R1, ChatGPT-5, ERNIE Bot 5.0, Qwen 3, dll.) secara mendalam mengubah cara penggunaan ponsel. Ponsel tidak lagi hanya alat untuk execute instruksi, tapi partner cerdas yang bisa memahami, menalar, dan berkreasi. Chip AI on-device (compute power NPU menembus 80 TOPS) membuat inference large model pertama kali berjalan lokal di ponsel, tanpa internet juga bisa pakai AI.

\booksubsection{Model besar AI (Large Language Model, LLM) adalah sistem AI yang dilatih dengan data teks masif, mampu memahami dan menghasilkan bahasa alami. Anda bisa membayangkannya sebagai “juara super” yang telah membaca hampir semua teks di internet, bisa menjawab pertanyaan, menulis artikel, menerjemahkan, memprogram, dan menganalisis data.}

Aplikasi Model Besar AI di Ponsel:

\booksubsection{Asisten percakapan}

\begin{itemize}
  \item \textbf{: App seperti DeepSeek, ChatGPT, Doubao, Kimi, dll., bisa bertanya, menerjemahkan, menulis copy, meringkas artikel kapan saja. Lebih cerdas dari mesin pencari tradisional—ia tidak hanya memberi Anda tautan, tetapi langsung memberi jawaban.}Penulisan dan kreasi
  \item \textbf{: Biarkan AI membantu Anda menulis email, meringkas notulen rapat, memoles artikel, menghasilkan copy kreatif.}Pembuatan dan pengeditan gambar
  \item \textbf{: AI menghapus orang yang lewat di foto dengan satu klik, memperluas batas gambar, mengubah gaya, bahkan menghasilkan gambar berdasarkan deskripsi teks.}Terjemahan real-time
  \item \textbf{Terjemahan real-time}Bimbingan belajar
  \item \textbf{: AI bisa menjelaskan konsep, memberikan soal latihan, mengoreksi pekerjaan—tutor pribadi yang selalu bersama Anda.}Model Besar Sisi Perangkat: Ponsel Menjalankan AI Sendiri
\end{itemize}

\booksubsection{Ponsel flagship terbaru sudah bisa menjalankan model besar AI kecil secara lokal (biasanya sekitar parameter 7B), perubahan yang dibawa:}

Tidak perlu jaringan

\begin{itemize}
  \item \textbf{: Bisa menggunakan fungsi AI bahkan saat offline—di pesawat atau di kereta bawah tanah.}Privasi lebih aman
  \item \textbf{: Percakapan dan data Anda tidak diunggah ke cloud, sepenuhnya diproses secara lokal di ponsel.}Respons lebih cepat
  \item \textbf{: Tanpa latensi jaringan, respons AI hampir instan.}Keterbatasan saat ini
  \item \textbf{: Kemampuan model sisi perangkat masih kalah dengan model besar cloud, tugas kompleks masih memerlukan koneksi internet. Namun seiring peningkatan daya komputasi AI chip, model sisi perangkat akan semakin kuat.}Integrasi AI Tingkat Sistem:
\end{itemize}

\booksubsection{Apple Intelligence}

\begin{itemize}
  \item \textbf{: Sistem AI yang diperkenalkan di iOS 18, mengintegrasikan alat tulis, magic wand gambar, Siri yang ditingkatkan, ringkasan notifikasi, dll. Memproses di sisi perangkat secara prioritas, tugas kompleks baru ke cloud.}Model Pangu Huawei
  \item \textbf{: Terintegrasi ke sistem HarmonyOS, asisten Xiaoyi mendapat kemampuan model besar, mendukung ringkasan dokumen, pengenalan layar cerdas, editing gambar AI, dll.}Surge AI Xiaomi
  \item \textbf{: Kemampuan AI tingkat sistem, mendukung pencarian gambar AI, ringkasan AI, subtitle real-time AI, percakapan model besar sisi perangkat, dll.}Samsung Galaxy AI
  \item \textbf{: Terjemahan panggilan real-time, ringkasan AI catatan, pembuatan dan pengeditan gambar, pencarian dengan lingkaran, dll.}5. Beberapa Skenario Aplikasi Tipikal dan Praktik
\end{itemize}

\booksection{Beberapa Skenario Aplikasi Tipikal dan Praktik}

Skenario 1: Terjemahan Real-time—Penerjemah Pribadi untuk Bepergian ke Luar Negeri

\booksubsection{Performa aktual}

\textbf{: Melihat bunga yang tidak dikenal di jalan, buka kamera, AI langsung memberi tahu “ini bunga hibiscus”. Mencari foto liburan tahun lalu di pantai di album, cari “laut”, AI otomatis menemukan semua foto yang mengandung laut, tanpa perlu mencari manual.}Prinsip di baliknya

\textbf{: Model computer vision (CV) melakukan vektorisasi (Embedding) pada konten gambar, lalu mencocokkan kemiripan dengan kueri teks. “Pencarian multimodal” ini adalah lompatan besar dalam manajemen album ponsel.}Skenario 2: Penghapusan AI—Menyelamatkan Foto Wisata yang “Dirusak” Orang Lewat

\booksubsection{Performa aktual}

\textbf{: Melihat bunga yang tidak dikenal di jalan, buka kamera, AI langsung memberi tahu “ini bunga hibiscus”. Mencari foto liburan tahun lalu di pantai di album, cari “laut”, AI otomatis menemukan semua foto yang mengandung laut, tanpa perlu mencari manual.}Prinsip di baliknya

\textbf{: Model computer vision (CV) melakukan vektorisasi (Embedding) pada konten gambar, lalu mencocokkan kemiripan dengan kueri teks. “Pencarian multimodal” ini adalah lompatan besar dalam manajemen album ponsel.}Skenario 3: Editing Foto AI—Satu Klik Jadi “Master Fotografi”

\booksubsection{Performa aktual}

\textbf{: Melihat bunga yang tidak dikenal di jalan, buka kamera, AI langsung memberi tahu “ini bunga hibiscus”. Mencari foto liburan tahun lalu di pantai di album, cari “laut”, AI otomatis menemukan semua foto yang mengandung laut, tanpa perlu mencari manual.}Prinsip di baliknya

\textbf{: Model computer vision (CV) melakukan vektorisasi (Embedding) pada konten gambar, lalu mencocokkan kemiripan dengan kueri teks. “Pencarian multimodal” ini adalah lompatan besar dalam manajemen album ponsel.}Skenario 4: Pemrosesan Dokumen—Pengganda Efisiensi untuk Pekerja Kantoran

\booksubsection{Performa aktual}

\textbf{: Melihat bunga yang tidak dikenal di jalan, buka kamera, AI langsung memberi tahu “ini bunga hibiscus”. Mencari foto liburan tahun lalu di pantai di album, cari “laut”, AI otomatis menemukan semua foto yang mengandung laut, tanpa perlu mencari manual.}Prinsip di baliknya

\textbf{: Model computer vision (CV) melakukan vektorisasi (Embedding) pada konten gambar, lalu mencocokkan kemiripan dengan kueri teks. “Pencarian multimodal” ini adalah lompatan besar dalam manajemen album ponsel.}Skenario 5: Pencarian Gambar dan Identifikasi Objek AI—“Mesin Pencari Visual” Anda

\booksubsection{Performa aktual}

\textbf{: Melihat bunga yang tidak dikenal di jalan, buka kamera, AI langsung memberi tahu “ini bunga hibiscus”. Mencari foto liburan tahun lalu di pantai di album, cari “laut”, AI otomatis menemukan semua foto yang mengandung laut, tanpa perlu mencari manual.}Prinsip di baliknya

\textbf{: Model computer vision (CV) melakukan vektorisasi (Embedding) pada konten gambar, lalu mencocokkan kemiripan dengan kueri teks. “Pencarian multimodal” ini adalah lompatan besar dalam manajemen album ponsel.}6. Optimasi Prompt: “Formula Komunikasi” agar AI Mendengarkan Perintah Anda

\booksection{Optimasi Prompt: “Formula Komunikasi” agar AI Mendengarkan Perintah Anda}

pemikiran komunikasi “Peran-Latar Belakang-Tugas-Batasan” yang paling praktis dan mudah dipelajari\textbf{, lengkap dengan banyak contoh.}Mengapa Perlu Kerangka?

\begin{tcbinfo}{Banyak orang menggunakan AI seperti menggunakan Baidu—memasukkan beberapa kata kunci. Tapi AI bukan mesin pencari, ia adalah model bahasa yang "memahami bacaan lalu menjawab".}
Beri lebih banyak konteks, jelaskan apa yang Anda inginkan, batasi agar tidak menyimpang\textbf{, ia bisa memberi hasil yang tepat. Kerangka adalah menerjemahkan “seperti apa jawaban yang baik” di pikiran Anda menjadi instruksi yang bisa dipahami AI.}1. Peran (Role): Beri Tahu AI Siapa Dirinya

\end{tcbinfo}

\booksubsection{Beri AI “identitas” spesifik, ia akan memanggil pengetahuan profesional dan cara ekspresi di bidang yang sesuai.}

Semakin spesifik pengaturan peran, semakin profesional outputnya\textbf{.}.

\booksubsection{Contoh Praktis:}

\begin{itemize}
  \item dokter TCM yang sabar, lembut, dan tidak pernah menertawakan pemula\textbf{." (Masalah teknis)}"Anda adalah
  \item dokter TCM yang sabar, lembut, dan tidak pernah menertawakan pemula\textbf{." (Masalah popularisasi sains)}"Anda adalah
  \item dokter TCM yang sabar, lembut, dan tidak pernah menertawakan pemula\textbf{." (Masalah penulisan)}"Anda adalah
  \item dokter TCM yang sabar, lembut, dan tidak pernah menertawakan pemula\textbf{." (Masalah kehidupan)}"Anda adalah
  \item dokter TCM yang sabar, lembut, dan tidak pernah menertawakan pemula\textbf{." (Masalah kesehatan)}Mengapa efektif?
\end{itemize}

\textbf{Mengapa efektif?}rentang pengetahuan\textbf{,},\textbf{, dan}“dalam kondisi apa”\textbf{AI. AI tanpa peran adalah “generalis yang serba bisa tapi tidak ahli”; dengan peran, ia bisa “masuk” ke identitas tersebut dan berpikir seperti profesional tersebut.}2. Latar Belakang (Background): Beri Tahu AI Situasi Anda

\booksubsection{Jelaskan sebelumnya “siapa Anda, situasi Anda saat ini, batasan Anda”, AI bisa memberikan jawaban yang}

benar-benar cocok untuk Anda\textbf{, bukan sekadar umum.}Contoh Praktis:

\booksubsection{Contoh Praktis:}

\begin{itemize}
  \item Xiaomi 14, dibeli setengah tahun lalu, baru-baru ini panas berlebih saat bermain game\textbf{." (Situasi perangkat)}"Saya adalah
  \item tante pensiunan berusia 60 tahun, mata tidak terlalu baik, tidak terlalu paham menggunakan ponsel pintar\textbf{." (Profil pengguna)}"Saya sedang
  \item mempersiapkan ujian CET-6 bahasa Inggris bulan Juni 2026, dasar biasa saja\textbf{." (Latar belakang belajar)}"Saya membuka
  \item kafe kecil di kota tier-3, omzet bulanan 30 ribu, mempekerjakan dua karyawan\textbf{." (Latar belakang bisnis)}"Saya
  \item baru saja melahirkan dan sedang masa pemulihan di rumah, kurang tidur parah\textbf{." (Tahap kehidupan)}Mengapa efektif?
\end{itemize}

\textbf{Mengapa efektif?}“untuk siapa melayani”\textbf{dan}“dalam kondisi apa”\textbf{. Soal yang sama, cara menjelaskan untuk siswa dan orang tua berbeda; saran untuk orang kota dan orang desa juga berbeda. Tanpa latar belakang, AI hanya bisa memberi jawaban “faktor persekutuan terbesar”, sering tidak mengena.}3. Tugas (Task): Jelaskan dengan Jelas Apa yang Harus Dilakukan AI

\booksubsection{Langkah ini gunakan}

awalan kata kerja\textbf{(tulis, daftar, bandingkan, terjemahkan, jelaskan, rangkum...), jelaskan tugas secara}spesifik dan dapat dieksekusi\textbf{. Tugas yang ambigu hanya akan menghasilkan jawaban yang ambigu.}Contoh Praktis:

\booksubsection{Contoh Praktis:}

\begin{itemize}
  \item tulis tiga copy video pendek untuk toko sarapan saya\textbf{, dan jelaskan alasannya." (Tindakan jelas)}"Tolong
  \item tulis tiga copy video pendek untuk toko sarapan saya\textbf{, pertahankan nada bisnis." (Tugas spesifik)}"Tolong
  \item tulis tiga copy video pendek untuk toko sarapan saya\textbf{, berikan tabel perbandingan kelebihan dan kekurangan." (Format jelas)}"Tolong
  \item tulis tiga copy video pendek untuk toko sarapan saya\textbf{." (Gaya jelas)}"Tolong
  \item tulis tiga copy video pendek untuk toko sarapan saya\textbf{, masing-masing untuk pekerja kantoran, pasangan, dan orang tua dengan anak}." (Segmentasi jelas)
\end{itemize}

\textbf{Mengapa efektif?}Bagian tugas adalah “pusat instruksi” tindakan AI.\textbf{Semakin jelas kata kerja, semakin spesifik permintaan, semakin AI tahu “apa yang harus disampaikan”}. “Lakukan sesuatu”, “buatkan”, “lihat-lihat” yang ambigu adalah tabu besar bagi AI—ia akan menebak maksud Anda, hasilnya sering menyimpang.

\booksubsection{Batasan (Constraints): Beri Tahu AI Apa yang Tidak Boleh Dilakukan}

“Batasan” terdengar seperti “kendala”, sebenarnya adalah\textbf{alat ampuh untuk membuat AI konvergen ke output yang paling cocok untuk Anda}. Dimensi umum meliputi: jumlah kata, format, nada, hal yang dihindari, hal yang harus dimasukkan.

\booksubsection{Contoh Praktis:}

\begin{itemize}
  \item "\textbf{Tidak lebih dari 200 kata, daftar dalam poin}." (Panjang + format)
  \item "\textbf{Nada harus ramah dan natural, seperti chat dengan teman}, jangan gunakan bahasa birokrasi seperti 'pertama, kedua, ketiga'." (Nada)
  \item "\textbf{Jangan gunakan istilah teknis}, jika menemui istilah profesional, jelaskan sekalian." (Batasan kata)
  \item "\textbf{Setiap paragraf tidak lebih dari 3 baris}, mudah dibaca di ponsel." (Tata letak)
  \item "\textbf{Harus memuat: skenario penggunaan, kelebihan dan kekurangan, kesesuaian dengan anggaran saya}." (Hal yang harus dimasukkan)
  \item "\textbf{Jangan berikan saran investasi spesifik}, hanya jelaskan prinsip." (Batasan)
\end{itemize}

\textbf{Mengapa efektif?}Tanpa batasan, AI mudah “bermain bebas”—entah terlalu akademis, terlalu panjang, atau menyimpang. Batasan yang jelas = “standar penyerahan” yang jelas, kedua pihak menyelaraskan ekspektasi, hasil bisa dikontrol.

\booksubsection{Evaluasi Kualitas Prompt: 5 Karakteristik Prompt yang Baik}

Menulis prompt yang baik bukan hal sekali jadi,\textbf{setelah selesai, periksa dulu lima hal ini}, baru kirim ke AI:

\begin{enumerate}
  \item \textbf{Spesifik bukan abstrak}: Bukan "tulis artikel", melainkan "tulis panduan pembelian ponsel 800 kata, 3 sub-bagian, untuk ibu baru pemula".
  \item \textbf{Termasuk latar belakang, bukan terisolasi}: Beri tahu AI "siapa saya, mengapa bertanya ini", ia bisa memberi jawaban "dibuat khusus untuk Anda".
  \item \textbf{Format output yang jelas}: Tabel, daftar, paragraf, atau JSON? Sebutkan sebelumnya.
  \item \textbf{Batasan yang tepat}: Jumlah kata, gaya, hal yang dihindari—pilih satu saja sudah memberi peningkatan yang nyata.
  \item \textbf{Dapat diverifikasi oleh AI}: Setelah bertanya, minta AI mengulangi permintaan Anda; jika tidak bisa menjelaskan dengan jelas, berarti tidak jelas.
\end{enumerate}

\booksubsection{Masalah Umum dan Solusi}

\begin{center}
\begin{tabularx}{\linewidth}{|X|X|}
\hline
\textbf{Gejala} & \textbf{Solusi} \tabularnewline
\hline
Jawaban terlalu umum, tidak praktis & Tambahkan latar belakang (siapa Anda/skenario apa) \tabularnewline
Jawaban terlalu akademis/terlalu resmi & Tambahkan batasan nada + daftar penolakan "bahasa birokrasi" \tabularnewline
Jawaban menyimpang/tidak menjawab pertanyaan & Ulangi tugas dengan kata kerja, dan minta "hanya jawab ini" \tabularnewline
Jawaban terlalu panjang, tidak bisa diselesaikan & Batasan jumlah kata yang keras (dalam 200 kata) \tabularnewline
AI bicara dengan sok tahu tapi salah & Tambahkan "jika tidak yakin, tandai 'saya tidak yakin'" + verifikasi fakta penting sendiri \tabularnewline
Versi pertama tidak memuaskan & Jangan tulis ulang, tambahkan "lebih ... lagi" di akhir prompt asli \tabularnewline
\hline
\end{tabularx}
\end{center}

\booksubsection{Tiga Kasus Skenario Lengkap}

\begin{tcbinfo}{Skenario 1: Asisten Kehidupan (Saran menu sarapan)}
\textbf{Prompt}:

Anda adalah blogger sarapan yang ahli dalam paduan nutrisi (\textbf{Peran}). Saya adalah ibu pekerja kantoran, setiap pagi hanya punya 15 menit untuk membuat sarapan, di rumah ada anak berusia 6 tahun dan mertua berusia 60 tahun (\textbf{Latar Belakang}). Tolong\textbf{daftar 7 menu sarapan tanpa pengulangan}(\textbf{Tugas}). Setiap sarapan harus: (1) Waktu persiapan tidak lebih dari 15 menit (2) Nutrisi seimbang (protein + karbohidrat + buah/sayur) (3) Cocok untuk semua usia (4) Bahan bisa dibeli di supermarket biasa (\textbf{Batasan}).

\textbf{Mengapa efektif}: Peran menentukan tingkat profesional, latar belakang memberi kesamaan tiga jenis orang, tugas menjelaskan hasil yang harus disampaikan, 4 batasan memastikan bisa dieksekusi.

\end{tcbinfo}

\begin{tcbinfo}{Skenario 2: Bantuan Belajar (Persiapan ujian CET-6 bahasa Inggris)}
\textbf{Prompt}:

Anda adalah\textbf{guru CET-6 bahasa Inggris yang berpengalaman}(Peran). Saya mahasiswa sarjana biasa,\textbf{skor CET-4 480, target CET-6 500, masih ada 3 bulan}(Latar belakang). Tolong\textbf{buatkan rencana persiapan 12 minggu yang bisa dieksekusi}(Tugas). Persyaratan: (1) Setiap minggu jelaskan alokasi waktu untuk mendengar, berbicara, membaca, menulis (2) Berikan fokus terobosan untuk setiap bulan (3) Rekomendasikan 3 buku materi review yang sesuai (jelaskan alasan dalam bahasa Mandarin) (4) Output dalam format tabel (Batasan).

\textbf{Mengapa efektif}: Peran menentukan latar belakang profesional, latar belakang memberi jangkar skor dan waktu, tugas menjelaskan “rencana bertahap”, batasan meminta “tabulasi” dan “bisa dieksekusi”.

\end{tcbinfo}

\begin{tcbinfo}{Skenario 3: Efisiensi Kerja (Penyusunan notulen rapat)}
\textbf{Prompt}:

Anda adalah\textbf{asisten administrasi senior}(Peran). Saya baru saja selesai rapat departemen 1 jam,\textbf{membahas rencana promosi pemasaran Q2}(Latar belakang).\textbf{Tolong susun catatan rapat berikut menjadi notulen terstruktur}(Tugas). Persyaratan: (1) Bagi menjadi tiga bagian “Keputusan / To-do / Risiko dan Keberatan” (2) To-do list tandai penanggung jawab dan tenggat waktu (jika tidak disebutkan, tulis “belum dikonfirmasi”) (3) Total kata dalam 600 kata (4) Output format markdown (Batasan).

\textbf{Mengapa efektif}: Peran menentukan perspektif administrasi, latar belakang memberi topik rapat, tugas menjelaskan output terstruktur, 4 batasan membuat hasil langsung bisa digunakan kembali.

\end{tcbinfo}

\textbf{Tingkat Lanjut: Jadikan AI “Asisten Pribadi” Anda}\par

Coba penggunaan lanjutan ini, rasakan kekuatan AI yang sebenarnya:

\begin{itemize}
  \item \textbf{Optimasi percakapan multi-putaran}: Saat tidak puas dengan jawaban AI, Anda bisa bilang "Ini terlalu profesional, tolong jelaskan sekali lagi dengan bahasa yang bisa dipahami siswa SD".
  \item \textbf{Chain of Thought (CoT)}: Minta AI "berpikir langkah demi langkah", kualitas jawaban biasanya lebih tinggi: "Tolong analisis soal matematika ini langkah demi langkah".
  \item \textbf{Kolaborasi multi-AI}: Gunakan A untuk menulis draf awal, B untuk memoles, C untuk menerjemahkan—kombinasi memberikan hasil yang lebih baik.
  \item \textbf{System prompt}: Preset identitas, preferensi, pantangan Anda di pengaturan AI, agar jawaban lebih sesuai kebutuhan Anda.
\end{itemize}

\booksection{Keterbatasan dan Tantangan AI: Pertahankan Penilaian yang Jernih}

AI memang kuat, tetapi bukan “serba bisa”. Memahami keterbatasannya dengan jernih akan membantu Anda menggunakannya dengan lebih baik:

\begin{itemize}
  \item \textbf{Masalah "halusinasi"}: AI bisa bicara sok tahu tapi salah (istilah profesional: Hallucination). Ia akan mengarang peristiwa sejarah yang tidak ada, kutipan yang salah, statistik palsu. Informasi penting harus diverifikasi secara manual.
  \item \textbf{Pemutusan pengetahuan}: Sebagian besar data pelatihan AI memiliki tanggal cutoff (misalnya awal 2024), ia tidak tahu apa yang terjadi setelah itu.
  \item \textbf{Kemampuan penalaran terbatas}: Penalaran multi-langkah yang kompleks, pembuktian matematika, jebakan logika—AI masih mudah salah.
  \item \textbf{Risiko privasi}: Pada layanan AI cloud, percakapan Anda mungkin digunakan untuk pelatihan model. Informasi sensitif sebaiknya prioritas menggunakan AI sisi perangkat.
  \item \textbf{Bias dan nilai}: AI mungkin membawa bias dari data pelatihan, perlu mempertahankan pemikiran independen.
\end{itemize}

\begin{tcbinfo}{Tip Penggunaan AI}
\begin{itemize}
  \item Konten yang dihasilkan AI mungkin tidak akurat (masalah "halusinasi"), informasi penting harus diverifikasi.
  \item Jangan memasukkan informasi privasi pribadi ke AI (nomor KTP, kata sandi, informasi bank, dll.).
  \item Manfaatkan prompt dengan baik: Semakin spesifik dan jelas deskripsi, semakin tinggi kualitas jawaban AI.
  \item AI adalah alat bukan otoritas, pertahankan pemikiran dan penilaian independen.
  \item Untuk keputusan kritis (medis, hukum, keuangan), konsultasikan dengan profesional daripada mengandalkan AI.
\end{itemize}

\end{tcbinfo}

\begin{tcbpractice}{Mari Tantang: Tiga Soal Praktik AI}
\begin{enumerate}
  \item \textbf{Tantangan 1: Instruksi kompleks}: Katakan pada asisten suara: "Besok pagi jam 7 ingatkan saya bawa payung, dan pindahkan kalender jam 3 sore hari ini ke jam 4." Lihat apakah ia bisa menyelesaikan operasi multi-langkah sekaligus.
  \item \textbf{Tantangan 2: Editing foto AI}: Ambil foto wisata dengan orang lewat, gunakan fungsi "Magic Eraser" atau "Penghapus AI", hapus orang lewat dengan bersih. Amati efek "imajinasi" AI pada latar belakang.
  \item \textbf{Tantangan 3: Optimasi prompt}: Pertama gunakan prompt sederhana "tulis panduan wisata" ke AI, lalu gunakan kerangka CRISPE untuk menulis prompt detail. Bandingkan perbedaan kualitas kedua jawaban.
\end{enumerate}

\end{tcbpractice}

\begin{tcbinfo}{Studi Kasus 2: Alur Lengkap Migrasi Data Ganti Ponsel}
\begin{itemize}
  \item \textbf{Skenario}: iPhone 13 Xiaomi 15 ganti ponsel lintas merek.
  \item \textbf{Langkah}: (1) Ponsel lama iCloud/WeChat backup bawaan (2) Logout Apple ID + matikan "Find" (3) Ponsel baru pakai App "Mi Mover" scan kode (4) Pilih item migrasi (5) WeChat terpisah pakai "Chat History Migration" bawaan (6) Verifikasi data kunci (kontak/album/WeChat) (7) Ponsel lama encrypted wipe.
  \item \textbf{Hasil yang Diharapkan}: Data zero loss, ponsel lama bisa aman dijual.
\end{itemize}

\end{tcbinfo}
